\worksheettitle{Домашняя работа: ответы}{\modulemark{M06-H} Версия учителя}

\begin{teacherbox}[Назначение]
  Домашняя работа проверяет построение максимума и минимума, положение нуля,
  повторяющиеся цифры, поразрядное сравнение и восстановление набора.
\end{teacherbox}

\section{Ответы}

\begin{teacherbox}[1. Составьте числа]
  Максимум \(9641\), минимум \(1469\); максимум \(8530\), минимум \(3058\);
  максимум \(8300\), минимум \(3008\); максимум \(7220\), минимум \(2027\).
  В каждой паре первым указан максимум, вторым \textemdash{} минимум.
\end{teacherbox}

\begin{teacherbox}[2. Сравнение]
  \begin{enumerate}
    \item \(5038<5308\). Первый различающийся разряд \textemdash{} сотни; \(0<3\).
    \item \(8610>8601\). Первый различающийся разряд \textemdash{} десятки; \(1>0\).
  \end{enumerate}
\end{teacherbox}

\begin{teacherbox}[3. Исправление]
  Минимум \(1058\). В числах \(1058\) и \(1508\) тысячи равны, а в разряде
  сотен \(0<5\), поэтому \(1058<1508\).
\end{teacherbox}

\begin{teacherbox}[4. Правило]
  По убыванию; нулём; все нули.
\end{teacherbox}

\newpage

\begin{teacherbox}[5. Несколько нулей и повторов]
  Максимум \(74400\), минимум \(40047\). В максимуме нули стоят справа как
  наименьшие цифры. В минимуме после первой ненулевой цифры 4 стоят оба нуля,
  затем оставшиеся цифры 4 и 7 идут по возрастанию.
\end{teacherbox}

\begin{teacherbox}[6. Восстановление]
  Даны цифры \(8,\ 6,\ 4,\ 0\). Минимум \(4068\).
\end{teacherbox}

\begin{teacherbox}[7. Две ошибки]
  Максимум \(9620\): в записи \(9622\) цифра 2 использована дважды, а цифра
  0 потеряна. Минимум \(2069\): запись \(0269\) начинается с нуля и обозначает
  трёхзначное число 269.
\end{teacherbox}

\begin{teacherbox}[8. Итог]
  Максимум \(97200\), минимум \(20079\). Принимается объяснение, в котором
  для максимума указан порядок по убыванию, а для минимума \textemdash{} первая
  ненулевая цифра 2, два нуля сразу после неё и порядок \(7<9\) в конце.
\end{teacherbox}

\begin{resultbox}[Критерии проверки]
  \begin{itemize}
    \item \textbf{Освоено:} все числа верны, использованы все цифры, объяснение
      связано с разрядами.
    \item \textbf{Точечная доработка:} одна перестановка при верном способе.
    \item \textbf{M06-R:} ведущий ноль, потеря цифры, случайное расположение
      нуля или отсутствие разрядного объяснения.
  \end{itemize}
\end{resultbox}
